% ============================================================
\chapter{Mod\`eles de S\'equences}
\label{ch:sequences}
% ============================================================

Le sac de mots jette l'ordre par la fen\^etre~: \og le chat mange la souris \fg et \og la souris mange le chat \fg ont la m\^eme repr\'esentation. Pour capturer le sens, il faut respecter la s\'equence. Les r\'eseaux de neurones r\'ecurrents (RNN), introduits par Jeffrey Elman en 1990, proposent une solution \'el\'egante~: un \'etat cach\'e qui se propage d'un mot au suivant, accumulant l'information contextuelle. Mais les RNN simples souffrent du probl\`eme du \emph{gradient \'evanescent}. La solution viendra des cellules LSTM (Hochreiter et Schmidhuber, 1997) et GRU (Cho et al., 2014), qui introduisent des m\'ecanismes de portes pour contr\^oler le flux d'information.

\begin{intuition}[title=\textbf{Mod\'eliser l'ordre des mots}]
Contrairement au mod\`ele sac-de-mots, les mod\`eles de s\'equences traitent
le texte comme une suite ordonn\'ee de tokens. Les r\'eseaux r\'ecurrents (RNN)
maintiennent un \'etat cach\'e qui encode l'historique de la s\'equence,
permettant de capturer les d\'ependances temporelles.
\end{intuition}

% ------------------------------------------------------------
\section{R\'eseaux de neurones r\'ecurrents (RNN)}
\label{sec:seq:rnn}
% ------------------------------------------------------------

\begin{definition}[RNN -- Elman Network]
Un r\'eseau r\'ecurrent simple (Elman, 1990) traite une s\'equence
$(\mathbf{x}_1, \ldots, \mathbf{x}_T)$ en maintenant un \'etat cach\'e
$\mathbf{h}_t \in \R^h$ :
\begin{align}
\mathbf{h}_t &= \tanh(\mathbf{W}_{hh} \mathbf{h}_{t-1} + \mathbf{W}_{xh} \mathbf{x}_t + \mathbf{b}_h) \\
\mathbf{y}_t &= \mathbf{W}_{hy} \mathbf{h}_t + \mathbf{b}_y
\end{align}
o\`u $\mathbf{W}_{hh} \in \R^{h \times h}$, $\mathbf{W}_{xh} \in \R^{h \times d}$,
$\mathbf{W}_{hy} \in \R^{o \times h}$ sont les matrices de poids.
\end{definition}

\begin{remarque}
L'\'etat cach\'e $\mathbf{h}_t$ agit comme une \og m\'emoire comprim\'ee \fg
de tout l'historique $(\mathbf{x}_1, \ldots, \mathbf{x}_t)$. En th\'eorie,
un RNN peut mod\'eliser des d\'ependances de longueur arbitraire ; en pratique,
cette capacit\'e est s\'ev\`erement limit\'ee.
\end{remarque}

% ------------------------------------------------------------
\section{Probl\`eme du gradient \'evanescent et explosif}
\label{sec:seq:gradient}
% ------------------------------------------------------------

\begin{theoreme}[Gradient \'evanescent]
Lors de la r\'etropropagation \`a travers le temps (BPTT), le gradient de la
perte $\mathcal{L}$ par rapport \`a $\mathbf{h}_k$ ($k \ll t$) satisfait :
\[
\frac{\partial \mathcal{L}_t}{\partial \mathbf{h}_k} = \frac{\partial \mathcal{L}_t}{\partial \mathbf{h}_t} \prod_{i=k+1}^{t} \frac{\partial \mathbf{h}_i}{\partial \mathbf{h}_{i-1}}
\]
o\`u chaque Jacobien $\frac{\partial \mathbf{h}_i}{\partial \mathbf{h}_{i-1}} = \text{diag}(\tanh'(\cdot)) \mathbf{W}_{hh}$.

Si $\norm{\mathbf{W}_{hh}} < 1 / \gamma$ (o\`u $\gamma = \max \abs{\tanh'(\cdot)}$),
le produit des Jacobiens d\'ecro\^it exponentiellement, rendant l'apprentissage
des d\'ependances longues impossible.
\end{theoreme}

\begin{proof}
On a $\norm{\frac{\partial \mathbf{h}_i}{\partial \mathbf{h}_{i-1}}} \leq \gamma \norm{\mathbf{W}_{hh}}$
o\`u $\gamma \leq 1$ pour $\tanh$. Donc :
\[
\norm{\prod_{i=k+1}^{t} \frac{\partial \mathbf{h}_i}{\partial \mathbf{h}_{i-1}}} \leq (\gamma \norm{\mathbf{W}_{hh}})^{t-k}
\]
Si $\gamma \norm{\mathbf{W}_{hh}} < 1$, cette quantit\'e tend vers $0$
exponentiellement vite quand $t - k \to \infty$.
\end{proof}

\begin{formules}[title=\textbf{Gradient clipping}]
Pour contrer le gradient explosif, on utilise le \emph{gradient clipping} :
\[
\hat{\mathbf{g}} = \begin{cases}
\mathbf{g} & \text{si } \norm{\mathbf{g}} \leq \theta \\
\theta \cdot \frac{\mathbf{g}}{\norm{\mathbf{g}}} & \text{sinon}
\end{cases}
\]
o\`u $\theta$ est le seuil de clipping (typiquement $\theta \in [1, 5]$).
\end{formules}

% ------------------------------------------------------------
\section{Long Short-Term Memory (LSTM)}
\label{sec:seq:lstm}
% ------------------------------------------------------------

\begin{definition}[LSTM]
Le r\'eseau LSTM (Hochreiter \& Schmidhuber, 1997) introduit une cellule
m\'emoire $\mathbf{c}_t$ et trois portes (gates) pour contr\^oler le flux
d'information :
\begin{align}
\mathbf{f}_t &= \sigma(\mathbf{W}_f [\mathbf{h}_{t-1}, \mathbf{x}_t] + \mathbf{b}_f) & \text{(porte d'oubli)} \\
\mathbf{i}_t &= \sigma(\mathbf{W}_i [\mathbf{h}_{t-1}, \mathbf{x}_t] + \mathbf{b}_i) & \text{(porte d'entr\'ee)} \\
\tilde{\mathbf{c}}_t &= \tanh(\mathbf{W}_c [\mathbf{h}_{t-1}, \mathbf{x}_t] + \mathbf{b}_c) & \text{(candidat)} \\
\mathbf{c}_t &= \mathbf{f}_t \odot \mathbf{c}_{t-1} + \mathbf{i}_t \odot \tilde{\mathbf{c}}_t & \text{(mise \`a jour)} \\
\mathbf{o}_t &= \sigma(\mathbf{W}_o [\mathbf{h}_{t-1}, \mathbf{x}_t] + \mathbf{b}_o) & \text{(porte de sortie)} \\
\mathbf{h}_t &= \mathbf{o}_t \odot \tanh(\mathbf{c}_t) & \text{(\'etat cach\'e)}
\end{align}
o\`u $\odot$ d\'enote le produit de Hadamard (\'el\'ement par \'el\'ement).
\end{definition}

\begin{proposition}[Gradient dans les LSTM]
Dans un LSTM, le gradient de la perte par rapport \`a la cellule m\'emoire
$\mathbf{c}_k$ contient un terme :
\[
\frac{\partial \mathbf{c}_t}{\partial \mathbf{c}_k} = \prod_{i=k+1}^{t} \left(\text{diag}(\mathbf{f}_i) + \ldots\right)
\]
La porte d'oubli $\mathbf{f}_i$ permet au gradient de se propager sans
att\'enuation quand $\mathbf{f}_i \approx \mathbf{1}$, r\'esolvant
partiellement le probl\`eme du gradient \'evanescent.
\end{proposition}

% ------------------------------------------------------------
\section{Gated Recurrent Unit (GRU)}
\label{sec:seq:gru}
% ------------------------------------------------------------

\begin{definition}[GRU]
Le GRU (Cho et al., 2014) simplifie le LSTM en fusionnant la cellule m\'emoire
et l'\'etat cach\'e, et en n'utilisant que deux portes :
\begin{align}
\mathbf{z}_t &= \sigma(\mathbf{W}_z [\mathbf{h}_{t-1}, \mathbf{x}_t] + \mathbf{b}_z) & \text{(porte de mise \`a jour)} \\
\mathbf{r}_t &= \sigma(\mathbf{W}_r [\mathbf{h}_{t-1}, \mathbf{x}_t] + \mathbf{b}_r) & \text{(porte de r\'einitialisation)} \\
\tilde{\mathbf{h}}_t &= \tanh(\mathbf{W}_h [\mathbf{r}_t \odot \mathbf{h}_{t-1}, \mathbf{x}_t] + \mathbf{b}_h) & \text{(candidat)} \\
\mathbf{h}_t &= (1 - \mathbf{z}_t) \odot \mathbf{h}_{t-1} + \mathbf{z}_t \odot \tilde{\mathbf{h}}_t & \text{(interpolation)}
\end{align}
\end{definition}

\begin{remarque}
Le GRU a moins de param\`etres que le LSTM ($3$ matrices de portes contre $4$)
et offre des performances comparables dans de nombreuses t\^aches. Le choix
entre LSTM et GRU est souvent empirique.
\end{remarque}

% ------------------------------------------------------------
\section{RNN bidirectionnels}
\label{sec:seq:bidirectional}
% ------------------------------------------------------------

\begin{definition}[BiRNN]
Un RNN bidirectionnel traite la s\'equence dans les deux directions et
concat\`ene les \'etats cach\'es :
\[
\mathbf{h}_t = [\overrightarrow{\mathbf{h}}_t ; \overleftarrow{\mathbf{h}}_t] \in \R^{2h}
\]
o\`u $\overrightarrow{\mathbf{h}}_t$ est l'\'etat de gauche \`a droite et
$\overleftarrow{\mathbf{h}}_t$ de droite \`a gauche.
\end{definition}

% ------------------------------------------------------------
\section{Architecture encodeur-d\'ecodeur (Seq2Seq)}
\label{sec:seq:seq2seq}
% ------------------------------------------------------------

\begin{definition}[Seq2Seq]
L'architecture encodeur-d\'ecodeur (Sutskever et al., 2014) utilise un RNN
encodeur pour compresser la s\'equence source en un vecteur de contexte
$\mathbf{c} = \mathbf{h}_T^{\text{enc}}$, et un RNN d\'ecodeur pour g\'en\'erer
la s\'equence cible :
\begin{align}
\text{Encodeur :} \quad \mathbf{h}_t^{\text{enc}} &= f_{\text{enc}}(\mathbf{x}_t, \mathbf{h}_{t-1}^{\text{enc}}) \\
\text{D\'ecodeur :} \quad \mathbf{h}_t^{\text{dec}} &= f_{\text{dec}}(\mathbf{y}_{t-1}, \mathbf{h}_{t-1}^{\text{dec}}) \\
P(y_t \mid y_{<t}, \mathbf{x}) &= \softmax(\mathbf{W}_o \mathbf{h}_t^{\text{dec}})
\end{align}
\end{definition}

\begin{attention}[title=\textbf{Goulot d'\'etranglement}]
Compresser toute la s\'equence source dans un seul vecteur $\mathbf{c}$ cr\'ee
un goulot d'\'etranglement informationnel. Ce probl\`eme est r\'esolu par le
m\'ecanisme d'attention (chapitre~\ref{ch:attention}).
\end{attention}

% ------------------------------------------------------------
\section{Impl\'ementation avec PyTorch}
\label{sec:seq:implementation}
% ------------------------------------------------------------

\begin{codeblock}[title=\textbf{LSTM pour la classification de texte}]
\begin{minted}{python}
import torch
import torch.nn as nn

class LSTMClassifier(nn.Module):
    def __init__(self, vocab_size, embed_dim, hidden_dim, num_classes,
                 num_layers=2, dropout=0.3, bidirectional=True):
        super().__init__()
        self.embedding = nn.Embedding(vocab_size, embed_dim, padding_idx=0)
        self.lstm = nn.LSTM(
            embed_dim, hidden_dim, num_layers=num_layers,
            batch_first=True, dropout=dropout,
            bidirectional=bidirectional
        )
        direction_factor = 2 if bidirectional else 1
        self.classifier = nn.Sequential(
            nn.Dropout(dropout),
            nn.Linear(hidden_dim * direction_factor, num_classes)
        )

    def forward(self, input_ids, lengths):
        embeds = self.embedding(input_ids)       # (B, T, d)
        packed = nn.utils.rnn.pack_padded_sequence(
            embeds, lengths.cpu(), batch_first=True, enforce_sorted=False
        )
        _, (hidden, _) = self.lstm(packed)       # hidden: (layers*dirs, B, h)
        # Concatenate last forward and backward hidden states
        hidden = torch.cat([hidden[-2], hidden[-1]], dim=1)  # (B, 2h)
        return self.classifier(hidden)           # (B, num_classes)

model = LSTMClassifier(
    vocab_size=30000, embed_dim=256,
    hidden_dim=128, num_classes=2
)
\end{minted}
\end{codeblock}

\begin{codeblock}[title=\textbf{Encodeur-d\'ecodeur Seq2Seq}]
\begin{minted}{python}
class Seq2SeqEncoder(nn.Module):
    def __init__(self, vocab_size, embed_dim, hidden_dim, num_layers=2):
        super().__init__()
        self.embedding = nn.Embedding(vocab_size, embed_dim)
        self.rnn = nn.GRU(embed_dim, hidden_dim, num_layers,
                          batch_first=True, bidirectional=True)
        self.fc = nn.Linear(hidden_dim * 2, hidden_dim)

    def forward(self, src):
        embedded = self.embedding(src)
        outputs, hidden = self.rnn(embedded)
        # Combine bidirectional hidden states
        hidden = torch.cat([hidden[-2], hidden[-1]], dim=1)
        hidden = torch.tanh(self.fc(hidden)).unsqueeze(0)
        return outputs, hidden

class Seq2SeqDecoder(nn.Module):
    def __init__(self, vocab_size, embed_dim, hidden_dim, num_layers=1):
        super().__init__()
        self.embedding = nn.Embedding(vocab_size, embed_dim)
        self.rnn = nn.GRU(embed_dim, hidden_dim, num_layers,
                          batch_first=True)
        self.fc_out = nn.Linear(hidden_dim, vocab_size)

    def forward(self, input_token, hidden):
        embedded = self.embedding(input_token).unsqueeze(1)
        output, hidden = self.rnn(embedded, hidden)
        prediction = self.fc_out(output.squeeze(1))
        return prediction, hidden
\end{minted}
\end{codeblock}

% ------------------------------------------------------------
\section{Teacher forcing et scheduled sampling}
\label{sec:seq:teacher}
% ------------------------------------------------------------

\begin{definition}[Teacher forcing]
Pendant l'entra\^inement, le \emph{teacher forcing} consiste \`a fournir au
d\'ecodeur le token de r\'ef\'erence $y_{t-1}$ plut\^ot que sa propre
pr\'ediction $\hat{y}_{t-1}$. Cela acc\'el\`ere la convergence mais cr\'ee
un d\'ecalage (\emph{exposure bias}) entre entra\^inement et inf\'erence.
\end{definition}

\begin{definition}[Scheduled sampling]
Le \emph{scheduled sampling} (Bengio et al., 2015) att\'enue l'exposure bias
en alternant al\'eatoirement entre teacher forcing et auto-r\'egression :
\[
\text{input}_t = \begin{cases}
y_{t-1} & \text{avec probabilit\'e } \epsilon_t \\
\hat{y}_{t-1} & \text{avec probabilit\'e } 1 - \epsilon_t
\end{cases}
\]
o\`u $\epsilon_t$ d\'ecro\^it progressivement de $1$ vers $0$.
\end{definition}

% ------------------------------------------------------------
\section{Exercices}
\label{sec:seq:exercices}
% ------------------------------------------------------------

\begin{exercice}[$\star$]
Comptez le nombre total de param\`etres d'un LSTM avec $d = 256$ (entr\'ee)
et $h = 512$ (\'etat cach\'e).
\end{exercice}

\begin{exercice}[$\star$]
Expliquez pourquoi un RNN bidirectionnel ne peut pas \^etre utilis\'e pour
la mod\'elisation de langue auto-r\'egressive.
\end{exercice}

\begin{exercice}[$\star\star$]
D\'emontrez que le gradient dans un RNN simple d\'ecro\^it exponentiellement
avec la distance temporelle sous les conditions du th\'eor\`eme du gradient
\'evanescent.
\end{exercice}

\begin{exercice}[$\star\star$]
Impl\'ementez un mod\`ele de langue caract\`ere par caract\`ere avec un LSTM
\`a 2 couches. Entra\^inez-le sur un corpus litt\'eraire et g\'en\'erez du
texte par \'echantillonnage.
\end{exercice}

\begin{exercice}[$\star\star\star$]
Comparez les performances d'un GRU et d'un LSTM sur une t\^ache de NER
(dataset CoNLL-2003). Analysez la vitesse de convergence et la capacit\'e
\`a capturer les d\'ependances longues.
\end{exercice}
